\fig{em6:line-element}{无损双导线的一个微段；串联电感和并联电容均按长度分配}{
\begin{tikzpicture}
\draw[world](0,0)--(7.1,0)
 (0,2.2)--(2,2.2) (3.6,2.2)--(7.1,2.2);
\foreach\j in {0,...,3}{\draw[world]({2+.4*\j},2.2)arc(180:0:.2);}
\node[above]at(2.8,2.55){$L'\dd x$};
\draw[world](4.5,2.2)--(4.5,1.3) (4.5,.9)--(4.5,0);
\draw[world](4.1,1.3)--(4.9,1.3) (4.1,.9)--(4.9,.9);
\node[right]at(5,1.1){$C'\dd x$};
\draw[flow](.3,2.6)--(1.5,2.6)node[midway,above]{$i(x,t)$};
\draw[flow](5.6,2.6)--(6.8,2.6)node[midway,above]{$i(x+\dd x,t)$};
\draw[dim](-.4,0)--(-.4,2.2)node[midway,left]{$u(x,t)$};
\draw[dim](7.5,0)--(7.5,2.2)node[midway,right]{$u(x+\dd x,t)$};
\draw[dim](0,-.7)--(7.1,-.7)node[midway,below]{$\dd x$};
\end{tikzpicture}}

图\ref{em6:line-element}中的串联电感造成相邻端电压的差，并联电容使流入、流出电流不同。取微段极限分别得到
\[
 \partial_xu=-L'\partial_t i,\qquad
 \partial_xi=-C'\partial_t u.
\]
再对第一式求 $x$ 偏导，并以第二式消去电流，
\[
 \partial_x^2u=-L'\partial_x\partial_t i
 =-L'\partial_t(-C'\partial_tu)=L'C'\partial_t^2u.
\]
令右行波 $u=f(\xi)$，$\xi=t-x/v$。链式法则给出
$\partial_xu=-f'(\xi)/v$、$\partial_x^2u=f''(\xi)/v^2$、
$\partial_t^2u=f''(\xi)$，代入波动方程得 $v^{-2}=L'C'$，
故 $v=1/\sqrt{L'C'}$。再令右行电流 $i=kf(\xi)$，第一式变为
$-f'/v=-L'kf'$，所以 $k=1/(vL')$，电压与电流之比为
$Z_0=vL'=\sqrt{L'/C'}$。这里选择零场为参考，未另加静态偏置。它来自图中电场、磁场储能的共同建立；图中没有电阻，$Z_0$仍然非零，所以不能把它理解为沿线发热电阻。
